\section{Heuristicas}
Es una función que estima que tan cerca esta un estado a su meta, guía para el camino con mayor promesa. Las heurísticas es que no sobrestimen los costos. Si dos estados tienen una evaluacipon heuristica similar -> preferiblemente examinar estado mas cercano a la raiz en grafo de estados.

\begin{definition}
Una \textbf{heurística} es una función que estima que tan cerca está un estado $n$ de la meta, diseñada especificamente para un problema de búsqueda particular.

\[
  h : \text{Estados} \;\longrightarrow\; \mathbb{R}_{\geq 0}
\]
En el contexto de búsqueda informada, $h(n)$ representa el \textbf{costo estimado} del camino mas barato desde el estado $n$ hasta la meta mas cercana. Por convención, $h(\text{meta}) = 0$
\end{definition}

En términos generales, una heurística es el estudio de metodos y reglas de descubrimiento e invención. En IA, nos interesa como \emph{aproximación informada} que:

\begin{itemize}
	\item usa conocimiento del dominio para guiar la busqueda.
	\item descarta/remueve posibilidades de solución poco prometedoras
	\item no garantiza exactitud, pero si eficiencia práctica
\end{itemize}

\begin{itemize}
		\item \textbf{distancia en linea recta}: útil en mapas de carreteras. $h(n) = $ distancia euclideana de $n$ a la meta.
		\item \textbf{distancia manhattan}: útil en grillas. $h(n) = \sum_i |x_i - x_{\text{meta}} | + |y_i -  y_{\text{meta}}|$
		\item \textbf{celdas fuera de lugar}: útil en el 8-puzzle. $h(n) = $ número de piezas no en su posición goal.
\end{itemize}

\subsection{Repaso algoritmos de búsqueda informada}
Los algoritmos estudiados en la clase anterior usan $h(n)$ de distintas formas:

\begin{center}
\begin{tabular}{lcc}
  \toprule
  \textbf{Algoritmo} & \textbf{Función de evaluación} & \textbf{Estrategia} \\
  \midrule
  Greedy Best-First & $f(n) = h(n)$          & Minimizar costo restante \\
  A*               & $f(n) = g(n) + h(n)$   & Minimizar costo total    \\
  A* con peso      & $f(n) = g(n) + w \cdot h(n)$ & Sesgar hacia la meta \\
  \bottomrule
\end{tabular}
\end{center}
 
donde $g(n)$ es el costo acumulado desde el inicio hasta $n$
(\emph{backward}), y $h(n)$ es la estimación del costo restante
(\emph{forward}).

\section{Admisibilidad, Monotonicidad, Informedness(Desigualdad triangular)}

\subsection{Admisibilidad}
\begin{definition}[Admisibilidad]
Una heurística $h$ es \textbf{admisible} (u \emph{optimista}) si nunca sobreestima el costo real para llegar a la meta.

\[
  \boxed{h(n) \;\leq\; \hstar{n}, \quad \forall n}
\]
donde  $\hstar{n}$ es el costo verdadero del camino optimo desde $n$ hasta la meta.
\end{definition}

\textbf{Consecuencia clave:} si $h$ es admisible, A* sobre arbol de busqueda es \emph{optimo}: nunca descarta la solución optima porque nunca penaliza en exceso ningun nodo.

\begin{example}[Admisibilidad en BFS]
BFS usa $f(n) = g(n)$ (equivalente a $h(n) = 0$ para todo $n$). Como
$0 \leq h^*(n)$ siempre, $h \equiv 0$ es admisible (aunque trivial e
ineficiente).
\end{example}

